\chapter{Related Work}\label{ch:related}

This chapter surveys existing work relevant to this thesis. It is organised thematically:
(i) trace-based interaction model extraction, (ii) scenario-driven and
functionality-centric analysis, (iii) behavioural pattern mining, (iv) static
analysis for architecture recovery (including the specific problem of
publish-subscribe dispatch), (v) hybrid static-dynamic approaches, and (vi) trace
visualisation for program comprehension.
For each area, the key techniques are described and their limitations with
respect to the problem addressed in this thesis are identified.
Representative studies for each area are summarised in \cref{tab:related-work-summary}.

The role of related work here is threefold. First, it grounds the approach in prior
techniques. Second, it identifies the gap that motivates this thesis. Third, it
establishes the baseline against which the results are compared.

\section{Trace-Based Interaction Model Extraction}

A significant body of work reconstructs \ac{UML} sequence or interaction
diagrams directly from execution traces.
Early work by \textcite{taniguchi:seqtrace} and \textcite{grati:interactive}
extracts sequence diagrams from Java execution traces. This demonstrates that
runtime information can be mapped onto object interactions.
\textcite{ziadi:dynamic} propose a fully dynamic reconstruction approach that
relies exclusively on traces with no static analysis support.
\textcite{hamou-lhadj:tracesummary} instead target scale. They summarise large
execution traces into compact \ac{UML} sequence diagrams by identifying and
merging recurring interaction patterns.
\textcite{baidada:petrinet} take a similar approach. They reconstruct higher-level
sequence diagrams by merging information from multiple execution traces and
modelling the result as Petri nets.
Briand \etal \cite{briand:seqdiag} use scenario-based testing to drive trace
collection and then synthesise sequence diagrams from the resulting traces.
This is scenario-specific but requires a scenario to be manually specified and
exercised up front.
\textcite{shang:loganalysis} show that even ordinary system logs (never
intended as execution traces) can be mined into architecture-level interaction
models. However, this comes at the cost of coarse granularity.

\subsection{Limitations}

All trace-based approaches share the same structural limitation: they can only
reconstruct interactions that were actually observed.
Missing events, low logging verbosity, and execution paths that were never
exercised during recording all show up as gaps in the reconstructed diagram.
In an industrial system with partial instrumentation, that is the common
case rather than the exception.

\section{Scenario-Driven and Functionality-Centric Analysis}

A separate line of work treats the \emph{scenario} itself as the organising
abstraction rather than the raw trace.
\textcite{egyed:scenario} introduces a scenario-driven approach to trace
dependency analysis. This approach explicitly links scenarios to execution footprints
and source-code artefacts. It demonstrates that scenarios can act as a unifying
concept across requirements, traces, and code.
Briand \etal \cite{briand:seqdiag}, already noted above, formalise the
mapping between runtime
events and interaction diagrams around execution scenarios in distributed
systems.
Despite recognising the value of scenario scoping, none of this work asks
whether static dependencies could be used to \emph{complete} a scenario-specific
interaction flow when the corresponding trace is incomplete or noisy.
The scenario bounds what is observed but does nothing to fill in what was not observed.

\section{Behavioural Pattern Mining and Model Inference}

A third line of research infers higher-level behavioural abstractions from
execution data rather than concrete interaction flows.
\textcite{lo:temporalrules} mine quantified temporal rules from execution
traces to extract recurring behavioural patterns across runs.
\textcite{gabel:javert} similarly mine general temporal properties expressed
as finite-state automata over method invocations. They do so fully automatically
from dynamic traces without requiring a pre-specified property to check against.
\textcite{walkinshaw:statemachines} reconstruct state-based behavioural
models (state machines) from software executions.
All three approaches infer generic behavioural structure purely from dynamic
traces. They use no static structural component and have no explicit grounding
in a specific functionality scenario.
The resulting models are rules, properties, or state machines rather than
concrete sequence diagrams. This makes them a poor fit for understanding
one particular execution flow.

\section{Static Analysis for Architecture Recovery}

Static analysis approaches extract architectural information directly from source
code without executing the system.
Formal call-graph construction techniques for \Cpp\ (including class hierarchy
analysis and points-to analysis) are surveyed in \textcite{grove:callgraph}.
These techniques produce over-approximate graphs that include every possible
interaction but cannot distinguish which interactions occur in a given scenario.

\textcite{ducasse:softvis} propose the \emph{class blueprint}, a polymetric
visualisation that overlays software metrics on a class's internal structure
(methods, attributes, and their invocation relationships) to support a first,
structural understanding of an unfamiliar class.
This line of work targets comprehension at the class level rather than
component-level interaction recovery and does not target \Cpp\ or
scenario-specific analysis.
Component-level dependency recovery is instead the focus of tools such as
Understand and Structure101, which produce rich static dependency graphs but
do not generate scenario-specific sequence diagrams.

\subsection{Handling Event-Driven Dispatch}

Static analysis of event-driven systems with publish-subscribe topologies is
particularly challenging. The subscription relationship between a
publisher and a handler is usually established indirectly (for example, via a
topic string or a registration callback) rather than through a direct call site.
\textcite{santos:roscompgraph} address this for \ac{ROS},
where inter-component communication is mediated by a topic broker.
Their approach statically extracts the ROS computation graph by matching
publisher and subscriber declarations by topic name to reconstruct the
communication topology.
This is a strategy closely related to the publish-subscribe topology
resolution developed in this thesis for CPSCore's IDC/Redis layer
(\cref{sec:idc}).

\section{Hybrid Static and Dynamic Analysis for Behaviour Recovery}

Hybrid approaches combining static and dynamic analysis aim to overcome the
completeness versus scenario-specificity trade-off directly.
\textcite{systa:reverse} combines static and dynamic views to support program
comprehension. This work shows that static structure can contextualise runtime
behaviour.
\textcite{hassan:changeprop} combine static and dynamic information to predict
change propagation. Their work demonstrates the benefit of integrating multiple
analysis sources for a different task.

\textcite{labiche:hybrid} propose a hybrid static-dynamic approach for
reverse-engineering scenario diagrams. They combine lightweight dynamic
instrumentation with static control-flow analysis to generate UML sequence
diagrams from Java systems.
Critically, static analysis in their approach is used to \emph{structure and
contextualise observed executions}. It is not used to infer interactions that
were never observed. This is the same gap identified in \cref{ch:introduction}.
\textcite{cornelissen:systematic} conduct a systematic review of dynamic
analysis for program comprehension and independently identify the lack of
integration with static information as a major open problem.
\textcite{kienle:reconstructing} survey the requirements that architecture
reconstruction tools must satisfy. They argue that effective tools need to
combine static structure with dynamic (trace) information to identify the
active subset of a system relevant to a given scenario.
Rather than presenting a concrete extraction technique, their contribution is
a requirements framework against which reconstruction tools can be evaluated.
It does not address \Cpp\ or publish-subscribe indirection specifically.
More recently, \textcite{abdelmalek:survey} surveys and categorises approaches
for reverse-engineering behavioural UML diagrams and identifies hybrid
static-dynamic techniques as a promising direction. However, the work
remains a categorisation exercise and stops short of a concrete method for
reconstructing scenario-specific models from incomplete traces.

Overall, hybrid analysis is widely acknowledged as beneficial. However,
existing approaches primarily enrich or organise already-observed runtime
behaviour. They do not systematically reason about incomplete traces or
reconstruct unobserved component interactions within a bounded execution
scenario.

\section{Trace Visualisation for Program Comprehension}

A further body of work improves program comprehension by presenting execution
traces at varying levels of abstraction, rather than by reconstructing a
formal model.
\textcite{cornelissen:controlled} and \textcite{cornelissen:systematic} study
trace visualisation techniques and their effectiveness for understanding
large execution traces.
\textcite{hamou-lhadj:tracesummary} propose trace summarisation techniques to
reduce noise and focus attention on relevant behaviour, and
\textcite{depauw:execpatterns} introduce visual representations of
object-oriented program execution aimed at scalability.
These techniques improve trace readability but, again, do not aim to
reconstruct a formal interaction model or to integrate static analysis to
complete missing behaviour.

\section{Summary and Positioning}

\Cref{tab:related-work-summary} summarises a representative selection of the
approaches reviewed above
along the dimensions most relevant to this thesis.

\begin{table}[htb]
\centering
\caption{Summary of representative related work along the five dimensions
  most relevant to this thesis.}
\label{tab:related-work-summary}
\small
\begin{tabularx}{\textwidth}{@{} l c c c c c @{}}
\toprule
\textbf{Approach} &
\textbf{Static} &
\textbf{Trace} &
\textbf{Hybrid} &
\textbf{Pub-sub} &
\textbf{SysML} \\
\midrule
Hamou-Lhadj \& Lethbridge \cite{hamou-lhadj:tracesummary}
  & $\times$ & \checkmark & $\times$ & $\times$ & $\times$ \\
Briand \etal \cite{briand:seqdiag}
  & $\times$ & \checkmark & $\times$ & $\times$ & $\times$ \\
Shang \etal \cite{shang:loganalysis}
  & $\times$ & \checkmark & $\times$ & $\times$ & $\times$ \\
Egyed \cite{egyed:scenario}
  & $\times$ & \checkmark & $\times$ & $\times$ & $\times$ \\
Lo \etal \cite{lo:temporalrules}
  & $\times$ & \checkmark & $\times$ & $\times$ & $\times$ \\
Gabel \& Su \cite{gabel:javert}
  & $\times$ & \checkmark & $\times$ & $\times$ & $\times$ \\
Ducasse \& Lanza \cite{ducasse:softvis}
  & \checkmark & $\times$ & $\times$ & $\times$ & $\times$ \\
Santos \etal \cite{santos:roscompgraph}
  & \checkmark & $\times$ & $\times$ & \checkmark & $\times$ \\
Syst\"{a} \cite{systa:reverse}
  & \checkmark & \checkmark & \checkmark & $\times$ & $\times$ \\
Labiche \etal \cite{labiche:hybrid}
  & \checkmark & \checkmark & \checkmark & $\times$ & $\times$ \\
Kienle \& M\"{u}ller \cite{kienle:reconstructing}
  & \checkmark & \checkmark & \checkmark & $\times$ & $\times$ \\
\midrule
\textbf{This thesis}
  & \checkmark & \checkmark & \checkmark & \checkmark & \checkmark \\
\bottomrule
\end{tabularx}
\end{table}

No existing approach addresses all five dimensions simultaneously.
In particular, three things are missing from prior work. First, none fuses static and trace information into a unified \acl{PG}.
Second, none applies dependence-based slicing to that graph for scenario
filtering. Third, none targets \sysml as the output formalism.
This gap constitutes the primary research contribution of this thesis.

